
% =====================================================================
\section{The Boundary--Thickness Theorem}
\label{sec:thickness}
% =====================================================================

We now prove the central structural fact: a bounded, finitely-resolved knower
cannot separate a region from its complement by a set of zero content. Every
separator it can realise has positive content, and that content is bounded
below by a constant determined by the resolution---the floor.

\subsection{Operational separators}

The topological boundary $\Bdy A$ can have measure zero (e.g.\ a half-space in
$\R^m$). The theorem is not about $\Bdy A$ in the ambient continuum; it is about
the separators a finite-resolution knower can actually realise, which must be
\emph{unions of cells}.

\begin{definition}[Realisable region and separator]
\label{def:realisable}
Fix a partition $\Part$ with resolution $\reso$. A region $A$ is
\emph{$\Part$-realisable} if it is a union of cells of $\Part$. A
\emph{$\Part$-separator} of a $\Part$-realisable $A$ is the set of cells of
$\Part$ that meet both $A$ and $\Phase\setminus A$ in the closure---formally,
\[
\Sigma_\Part(A)\;=\;\bigcup\bigl\{\cell\in\Part:\ \cl(\cell)\cap\cl(A)\neq\varnothing
\ \text{and}\ \cl(\cell)\cap\cl(\Phase\setminus A)\neq\varnothing\bigr\}.
\]
Because a finite-resolution knower can only point to cells, $\Sigma_\Part(A)$ is
the thinnest separator it can name.
\end{definition}

\begin{remark}
The shift from $\Bdy A$ to $\Sigma_\Part(A)$ is the whole content of
``operational.'' A continuum knower with unlimited resolution might use $\Bdy A$
of measure zero; a knower bound by \ref{bps:res} cannot, because it cannot name
sub-cell sets. The theorem quantifies the price of that limitation.
\end{remark}

\subsection{Non-triviality of the separator}

\begin{lemma}[The realisable separator is non-empty]
\label{lem:nonempty-sep}
Let $A$ be $\Part$-realisable with $\mu(A)>0$ and $\mu(\Phase\setminus A)>0$, and
suppose $\Phase$ is connected. Then $\Sigma_\Part(A)\neq\varnothing$.
\end{lemma}

\begin{proof}
Suppose $\Sigma_\Part(A)=\varnothing$. Then no cell meets the closures of both
$A$ and $\Phase\setminus A$; hence $\cl(A)$ and $\cl(\Phase\setminus A)$ are
disjoint unions of whole cells whose closures do not touch. Writing
$U=\interior(\cl A)$ and $V=\interior(\cl(\Phase\setminus A))$, the sets $U,V$
are non-empty (each contains a cell of positive measure), open, disjoint, and
cover $\Phase$ up to a $\mu$-null set; if their closures do not meet, $U\sqcup V$
is a separation of $\Phase$ into two non-empty clopen pieces, contradicting
connectedness. Hence $\Sigma_\Part(A)\neq\varnothing$.
\end{proof}

\begin{remark}[Connectedness]
\Cref{lem:nonempty-sep} uses connectedness of $\Phase$. If $\Phase$ is
disconnected, a region can coincide with a component and have empty separator;
but then $A$ was not genuinely \emph{divided} from its complement---it was given.
The theorem concerns acts of division \emph{within} a connected whole, which is
the case of interest (an object carved out of one world). We assume connectedness
hereafter and note it where used.
\end{remark}

\subsection{The thickness lower bound}

\begin{definition}[Boundary thickness]
\label{def:thickness}
The \emph{boundary thickness} of a $\Part$-realisable region $A$ is the measure
of its realisable separator,
\[
\thick_\Part(A)\;:=\;\mu\bigl(\Sigma_\Part(A)\bigr).
\]
\end{definition}

\begin{theorem}[Boundary--Thickness Theorem]
\label{thm:thickness}
Let $(\Phase,d,\mu)$ be a connected bounded phase space satisfying
\Cref{ax:bps}, and let $\Part$ be a realisable partition of resolution $\reso$
with minimum cell measure $\mu_{\min}>0$. Then for every $\Part$-realisable
region $A$ with $\mu(A)>0$ and $\mu(\Phase\setminus A)>0$,
\[
\boxed{\ \thick_\Part(A)\ \ge\ \mu_{\min}\ >\ 0.\ }
\]
In particular no realisable separator is sharp: $\mu(\Sigma_\Part(A))=0$ is
impossible.
\end{theorem}

\begin{proof}
By \Cref{lem:nonempty-sep}, $\Sigma_\Part(A)$ contains at least one cell
$\cell^\ast\in\Part$. By \ref{bps:res} and \Cref{lem:finite-cells} every cell
has $\mu(\cell)\ge\mu_{\min}>0$. Since $\Sigma_\Part(A)$ is a union of cells and
contains $\cell^\ast$,
\[
\thick_\Part(A)=\mu(\Sigma_\Part(A))\ \ge\ \mu(\cell^\ast)\ \ge\ \mu_{\min}\ >0 .
\]
A sharp separator would have $\mu(\Sigma_\Part(A))=0<\mu_{\min}$, a
contradiction.
\end{proof}

\begin{definition}[Resolution floor]
\label{def:floor}
The \emph{resolution floor} of a knower realising partition $\Part$ is
\[
\thick\ :=\ \inf_{A}\ \thick_\Part(A)\ \ge\ \mu_{\min}\ >\ 0,
\]
the infimum over all $\Part$-realisable regions of their boundary thickness. By
\Cref{thm:thickness} the floor is strictly positive.
\end{definition}

\begin{corollary}[No zero-width cut]
\label{cor:no-zero-cut}
Under \Cref{ax:bps}, a bounded finite-resolution knower cannot divide a
connected whole into a region and its complement by a separator of zero content.
Every division it performs deposits content $\ge\thick$ in the separator.
\end{corollary}

\begin{proof}
Immediate from \Cref{thm:thickness} and \Cref{def:floor}.
\end{proof}

\subsection{Stability and the role of boundedness}

The next two results show the bound is robust and that boundedness is essential,
not decorative.

\begin{proposition}[Refinement lowers but never abolishes the floor]
\label{prop:refine}
Let $\Part'\preceq\Part$ be a refinement with $\mu'_{\min}\le\mu_{\min}$. Then
$\thick'\le\thick$ for the refined floor, and still $\thick'\ge\mu'_{\min}>0$.
The floor decreases monotonically under refinement but is bounded below by the
finest realisable cell and never reaches zero at any finite stage.
\end{proposition}

\begin{proof}
A finer partition can only shrink the separating cells, so
$\thick'_{\Part'}(A)\le\thick_\Part(A)$ for each $A$, giving $\thick'\le\thick$.
By \Cref{thm:thickness} applied to $\Part'$, $\thick'\ge\mu'_{\min}>0$. By
\ref{bps:res} every finite stage carries a positive resolution and hence a
positive $\mu'_{\min}$; only in the (non-realisable) limit $\reso\to0$ would the
bound vanish, and that limit is excluded for any embedded knower.
\end{proof}

\begin{proposition}[Boundedness is necessary]
\label{prop:bdd-necessary}
The conclusion of \Cref{thm:thickness} fails without \ref{bps:bdd}. On an
unbounded space $\R^m$ with Lebesgue measure there exist regions (e.g.\
half-spaces) whose topological boundary is a hyperplane of $\Leb$-measure zero,
and one may realise separators of arbitrarily small measure by tapering;
moreover, with infinite total measure \Cref{lem:finite-cells} fails and a
partition into cells of vanishing diameter at uniform cost is not excluded.
\end{proposition}

\begin{proof}
Take $A=\{x_1\le0\}\subset\R^m$. Then $\Bdy A=\{x_1=0\}$ has
$\Leb(\Bdy A)=0$, so the topological separator is sharp. The finite-resolution
obstruction of \Cref{thm:thickness} relied on \Cref{lem:finite-cells}, whose
proof used $\mu(\Phase)<\infty$; with $\mu(\R^m)=\infty$ the measures of
infinitely many cells can sum to $\infty$ without forcing
$\inf_i\mu(\cell_i)=0$, so the uniform lower bound $\mu_{\min}>0$ need not hold.
Hence the floor argument does not apply, witnessing the necessity of
\ref{bps:bdd}.
\end{proof}

\begin{remark}[What the theorem does and does not say]
\label{rem:thickness-scope}
\Cref{thm:thickness} does not assert that the ambient continuum has thick
boundaries; a half-space in $\R^m$ has a measure-zero topological boundary. It
asserts that a \emph{bounded finite-resolution knower}, who can name only unions
of positive-measure cells, cannot realise that measure-zero cut: its sharpest
nameable separator has measure $\ge\mu_{\min}$. The thickness is a property of
the knower-in-the-world, jointly with boundedness, not of the world in the
abstract. This is the precise sense in which ``items are bounded'': the divisions
by which items are individuated are bounded below in content.
\end{remark}
